\chapter{Memory, mutability and references}
\label{chap:mutability}

Chapter~\ref{chap:collections} gave a rule without its reason: a collection
whose elements change is declared \token{dmut}, and \token{mut} is not enough.
It also left a warning without its explanation: a grid made by repeating one
row has rows that change together. Both come from the same fact. Some values
are not stored in their variable, and two variables can then refer to the same
elements. A change made through one of them is seen through the other.

This chapter shows where a program keeps its values, and what happens to them
when a variable is assigned to another, or passed to a function. It then
presents the tools Ymir gives to control it: the two levels of \token{mut} and
\token{dmut}, \token{alias} to share elements that can change, \token{copy} and
\token{dcopy} to stop sharing them, and references, with which a function
changes a variable of its caller. Loops and patterns take the same keywords, to
change the elements they go through or take apart. The chapter ends with a
program that changes a grid in place, generation after generation: the game of
life.

Part II covers the same ground as a specification, in
Chapter~\ref{chap:memory}.

\minitoc%

\vfill\pagebreak
\section{Where values are kept}
\label{sec:mutability:stack_heap}

A running program keeps its values in the memory of the machine, a long row of
bytes, each with a number called its \textit{address}. The program uses two
regions of it for the values of its variables.

\begin{itemize}
\item \textit{The stack.} Each call of a function takes a zone of the stack for
  its parameters and its variables, called its \textit{frame}. The frame is
  taken when the call starts and given back when it returns, so the variables
  of a function exist only during its call. A call made from inside the
  function gets a frame of its own, next to that of its caller: this is the
  call stack of Section~\ref{sec:functions:infinite_recursion}, which
  overflows when too many calls wait at the same time.
\item \textit{The heap.} The program can also reserve a block of memory
  anywhere in the heap, of any size, when it runs. Chapter~\ref{chap:collections}
  called that an allocation. A block does not belong to a call: it stays as long
  as the program needs it, whichever functions return in the meantime.
\end{itemize}

The next listing declares an array and a slice, and
Figure~\ref{fig:mutability:array_slice} shows where their elements are.

\lstinputlisting[style=coloredverbatim, caption={An array and a slice, \textit{examples/mutability/where.yr}}, label=lst:mutability:where]{examples/mutability/where.yr}
\vspace{-10pt}%
\begin{lstlisting}[style=bashVerb]
[31, 28, 31] [2, 3, 5, 7]
\end{lstlisting}

\begin{figure}[h]
  \centering
  \begin{adjustbox}{max width=\linewidth}
    \begin{tikzpicture}
      \foreach \value [count=\i from 0] in {31, 28, 31} {
        \node[memcell] (d\i) at (\i * 1.0, 0) {\value};
      }
      \node[cflab, left=3mm of d0] (dn) {\token{days}};
      \node[memcell] (len) at (0, -1.2) {4};
      \node[memaddr, right=0mm of len] (ptr) {};
      \node[cflab, left=3mm of len] (pn) {\token{primes}};
      \node[memnote, below=0.5mm of len] {length};
      \node[memnote, below=0.5mm of ptr] {address};
      \node[memframe, fit=(dn)(d2)(pn)(ptr), inner ysep=4mm] (main) {};
      \node[memtitle] at ([xshift=2mm]main.north west) {main};

      \foreach \value [count=\i from 0] in {2, 3, 5, 7} {
        \node[memcell] (h\i) at (5.2 + \i * 1.0, -1.2) {\value};
      }
      \draw[mempoint] (ptr.center) -- (h0.west);

      \node[memregion] (stack) at ($(main.north) + (0, 5mm)$) {stack};
      \node[memregion] at ({$(h0)!0.5!(h3)$} |- stack.base) {heap};
      \draw[memsep] (4.3, 0.9) -- (4.3, -2.4);
    \end{tikzpicture}
  \end{adjustbox}
  \caption{\label{fig:mutability:array_slice} The array \token{days} and the
    slice \token{primes} of Listing~\ref{lst:mutability:where}. The elements of
    the array are in the variable, in the frame of \token{main} on the stack.
    The slice holds a length and an address, and its elements are in a block
    of the heap.}
\end{figure}


\subsection{Values stored in the variable}

The three elements of \token{days} are in the variable itself, one after the
other, in the frame of \token{main}. The same goes for every type met so far
except slices and maps: an integer, a float, a character or a boolean take a few
bytes of the frame, and a tuple or an array take the bytes of all their fields
or elements. The compiler knows that size, since the length of an array is part
of its type, and reserves the frame accordingly.

\subsection{Values stored elsewhere}

A slice cannot be stored that way: its length is known only when the program
runs, and it can grow. The variable \token{primes} holds two numbers of a fixed
size instead: the length of the slice, and the address of its first element.
The elements are in a block of the heap, one after the other, as those of an
array are. Reading \token{primes[2]} follows the address, and counts two
elements from there. (A slice keeps a third number, used when it grows, which
the figures of this chapter leave out.)

A value whose elements lie outside of it is said to \textit{borrow} them: the
elements are not part of the value, the value only says where they are. Slices,
strings, which are slices of characters, and maps borrow their elements. A map
is laid out differently from a slice, but the principle is the same: its
variable holds the address of a structure of the heap, where the entries are.

The empty slice \token{[]} has a length of 0 and no element to designate, so it
needs no block, and no \token{copy}. A map refers to its structure even when it
has no entry, and \token{copy []} is what allocates that structure, which is
why an empty map cannot do without it (Section~\ref{sec:collections:maps}).

This answers a question of Chapter~\ref{chap:collections}. \token{[2, 3, 5, 7]}
is an array: its four elements are the value itself, and would be kept in the
frame. \token{copy} reserves a block of the heap, copies the four elements into
it, and gives a slice that borrows them. The block outlives the call that made
it, which is why a function can return a slice that it created.

A block of the heap is never given back by the program itself. The
\textit{garbage collector}, a part of every Ymir program, finds the blocks that
no value refers to any more, and reuses their memory. A program never has to
free a block, and never reads one that was freed.

\notebox{A string literal such as \token{"hello"} is a slice too, but its
  characters are not in the heap: the compiler writes them into the program
  itself, in a region that cannot change. \token{copy "hello"} copies them into
  a block of the heap, where they can change.}

\vfill\pagebreak
\section{Two names for the same elements}
\label{sec:mutability:sharing}

Assigning a variable to another copies its value. What that copies depends on
where the value is. The next listing assigns an array, then a slice, and
changes the first variable afterwards.

\lstinputlisting[style=coloredverbatim, caption={Assigning an array and a slice, \textit{examples/mutability/share.yr}}, label=lst:mutability:share]{examples/mutability/share.yr}
\vspace{-10pt}%
\begin{lstlisting}[style=bashVerb]
[31, 29, 31] [31, 28, 31]
[7, 3, 5] [7, 3, 5]
\end{lstlisting}

The array behaves as an integer would. \token{let before = days;} on line~5
copies the three elements into \token{before}, and line~6 changes
\token{days[1]} only: \token{before} keeps the 28 it received.

The slice does not. \token{let seen = primes;} on line~10 copies what
\token{primes} holds, which is a length and an address. \token{seen} designates
the same elements as \token{primes}, and when line~11 changes the first of them,
\token{seen[0]} is 7 as well. Figure~\ref{fig:mutability:shared} shows both
cases.

\begin{figure}[h]
  \centering
  \begin{adjustbox}{max width=\linewidth}
    \begin{tikzpicture}
      \foreach \value [count=\i from 0] in {31, 29, 31} {
        \ifnum\i=1
          \node[memhit] (d\i) at (\i * 1.0, 0) {\value};
        \else
          \node[memcell] (d\i) at (\i * 1.0, 0) {\value};
        \fi
      }
      \node[cflab, left=3mm of d0] (dn) {\token{days}};
      \foreach \value [count=\i from 0] in {31, 28, 31} {
        \node[memcell] (b\i) at (\i * 1.0, -1.0) {\value};
      }
      \node[cflab, left=3mm of b0] {\token{before}};

      \node[memcell] (plen) at (0, -2.2) {3};
      \node[memaddr, right=0mm of plen] (pptr) {};
      \node[cflab, left=3mm of plen] (pn) {\token{primes}};
      \node[memcell] (slen) at (0, -3.2) {3};
      \node[memaddr, right=0mm of slen] (sptr) {};
      \node[cflab, left=3mm of slen] (sn) {\token{seen}};
      \node[memframe, fit=(dn)(d2)(sn)(sptr)] (main) {};
      \node[memtitle] at ([xshift=2mm]main.north west) {main};

      \foreach \value [count=\i from 0] in {7, 3, 5} {
        \ifnum\i=0
          \node[memhit] (h\i) at (5.2 + \i * 1.0, -2.7) {\value};
        \else
          \node[memcell] (h\i) at (5.2 + \i * 1.0, -2.7) {\value};
        \fi
      }
      \draw[mempoint] (pptr.center) to[out=0, in=180] (h0.west);
      \draw[mempoint] (sptr.center) to[out=0, in=180] (h0.west);

      \node[memregion] (stack) at ($(main.north) + (0, 5mm)$) {stack};
      \node[memregion] at (h1 |- stack.base) {heap};
      \draw[memsep] (4.3, 0.9) -- (4.3, -4.0);
    \end{tikzpicture}
  \end{adjustbox}
  \caption{\label{fig:mutability:shared} The variables of
    Listing~\ref{lst:mutability:share} at its end. \token{before} is a copy of
    the elements of \token{days}, and did not follow the change of
    \token{days[1]}. \token{seen} is a copy of the length and the address of
    \token{primes}: both designate the same elements, and \token{primes[0] = 7;}
    changed what both hold.}
\end{figure}


Sharing is what makes slices cheap to pass around: assigning one or passing it
to a function copies two numbers, whatever its length. The functions
\token{sum} and \token{average} of Listing~\ref{lst:collections:average}
received the slice of their caller that way, and did not copy its elements. A
part of a slice, \token{s[1 .. 4]}, shares the elements of \token{s} in the same
way (Section~\ref{sec:collections:slices}).

As long as no one changes the elements, sharing them cannot be seen: two
variables that designate the same elements \token{[2, 3, 5]} and two variables
that hold a copy each print the same thing. It shows once one of them changes
an element, as \token{primes} does on line~11. \token{seen} is a constant, and
cannot change the elements itself, but it sees the changes made through
\token{primes}. The next sections explain which variables may change shared
elements, and how a program says that it expects them to.

\subsection{Giving a variable a new slice}

Changing an element and giving the variable a new value are two different
things. Changing an element writes into the block that the slice designates, and
every variable that designates it sees the change. Giving the variable a new
value replaces the length and the address that it holds, and the others are not
affected.

\begin{lstlisting}[style=coloredverbatim]
let mut primes = copy [2, 3, 5];
let seen = primes;
primes = copy [11, 13];
println(primes, " ", seen);
\end{lstlisting}
\vspace{-10pt}%
\begin{lstlisting}[style=bashVerb]
[11, 13] [2, 3, 5]
\end{lstlisting}

After the third line, \token{primes} designates a new block of two elements,
and \token{seen} still designates the first one.

\tokennolst{\mtildeascii{}=} (Section~\ref{sec:collections:slices}) also gives
the variable a new slice, but not always a new block. When the block of the
slice has room after its last element, the elements appended are written there,
and the variable designates the same block, one element longer. A variable that
shared the elements before still shares the first ones.

\begin{lstlisting}[style=coloredverbatim]
let dmut numbers = copy [1, 2, 3];
let seen = numbers;
numbers ~= [4];
numbers[0] = 9;
println(numbers, " ", seen);
\end{lstlisting}
\vspace{-10pt}%
\begin{lstlisting}[style=bashVerb]
[9, 2, 3, 4] [9, 2, 3]
\end{lstlisting}

Whether the block has room depends on how it was allocated, and a program
should not count on either outcome. When two variables must not share anything,
one of them is made with \token{copy}, which Section~\ref{sec:mutability:copies}
presents.

\vfill\pagebreak
\section{Two levels of mutability}
\label{sec:mutability:levels}

Section~\ref{sec:collections:dmut} gave a table of what may change in a
collection, depending on its declaration. The memory of the previous sections
explains it. A slice has two parts: the length and the address held by the
variable, and the elements, in their block. Each part has its own mutability.

\begin{center}
  \begin{adjustbox}{max width=\linewidth}
    \begin{tabular}{|l|l|l|l|}
      \hline
      Declaration & Type & The variable & The elements\\[0pt]
      \hline
      \hline
      \token{let s = copy [1, 2];} & \token{[i32]} & constant & constant\\[0pt]
      \token{let mut s = copy [1, 2];} & \token{mut [i32]} & mutable & constant\\[0pt]
      \token{let dmut s = copy [1, 2];} & \token{mut [mut i32]} & mutable & mutable\\[0pt]
      \hline
    \end{tabular}
  \end{adjustbox}
\end{center}

A \token{mut} in a type makes mutable the level that follows it. In
\token{mut [mut i32]}, the first \token{mut} applies to the slice itself, the
variable, which can then be given another slice, and the second to its elements,
the \token{i32} of the block. \token{let mut} makes the first level mutable,
and \token{let dmut}, for \textit{deeply mutable}, makes every level mutable.
The messages of the compiler write types in full, levels included: the
\token{mut [mut i32 ; 4us]} of \errmsg{E4095}{incompatible types [i32] and mut
  [mut i32 ; 4us]} (Section~\ref{sec:collections:copy}) is an array literal,
which may be changed at both levels since nothing else refers to it yet.

A slice of slices has three levels: the variable, the rows, and the elements of
the rows. \token{dmut grid: [[i32]]} is a \token{mut [mut [mut i32]]}, and all
three can change: \token{grid = ...}, \token{grid[1] = ...} and
\token{grid[1][2] = ...}.

An array has levels too, although its elements are in the variable: the type
of \token{let dmut counts = [0 ; 6];} is \token{mut [mut i32 ; 6]}, and
\token{let mut} alone does not let its elements change. An integer has only one
level, and \token{let dmut n = 1;} means the same as \token{let mut n = 1;}.

\subsection{Why two levels}

A variable that can take a new value affects no other variable. Giving
\token{primes} a new slice, as in Section~\ref{sec:mutability:sharing}, leaves
\token{seen} as it was. Changing the elements is different: they may be shared,
and every variable that shares them sees the change. The two levels keep these
two kinds of change apart. \token{let mut} allows only the first, which is
always local to the variable, and \token{let dmut} is written where the second
is wanted as well.

This is the reason for the rule of Chapter~\ref{chap:collections}: a program
that changes elements must say so at the declaration of the collection, and a
reader who meets a variable declared \token{let} or \token{let mut} knows that
nothing it designates changes through it.

\vfill\pagebreak
\section{Sharing what can change: \texttt{alias}}
\label{sec:mutability:alias}

In Listing~\ref{lst:mutability:share}, \token{seen} shares the elements of
\token{primes} but cannot change them. A second variable that may change them
too is refused, if it is declared as \token{seen} was.

\begin{lstlisting}[style=coloredverbatimError, escapechar=@, caption=Two variables that may change the same elements]
use std::io;

fn main() {
  let dmut scores = copy [12, 15, 9];
  let dmut same = @\hb{scores}@; // error
  same[2] = 10;
  println(scores, " ", same);
}
\end{lstlisting}

The compiler says \errmsg{E4045}{discarding the constant property is
  prohibited}, and the note under it gives the reason: \textit{implicit alias of
  type mut [mut i32] is prohibited, it implicitly gives up on borrowings of
  mutable values}. An \textit{alias} is a second name for the same elements.
The declaration would create one silently, and the line \token{same[2] = 10;}
would then change \token{scores}, which it does not name.

A reader of a long function cannot keep in mind every pair of variables that
share elements. Ymir therefore asks for such a pair to be created in writing,
with the keyword \token{alias} in front of the value shared.

\lstinputlisting[style=coloredverbatim, caption={A second name for the same elements, \textit{examples/mutability/alias.yr}}, label=lst:mutability:alias]{examples/mutability/alias.yr}
\vspace{-10pt}%
\begin{lstlisting}[style=bashVerb]
[12, 15, 10] [12, 15, 10]
\end{lstlisting}

\token{alias scores}, on line~5, gives the length and the address of
\token{scores}, as a plain assignment would. The keyword changes nothing to what
is copied; it records that the program means the two variables to share
elements that both may change. A reader who looks for the places where
\token{scores} may change without being named finds them by looking for
\token{alias scores}.

\subsection{Only where it is needed}

\token{alias} is required where a value whose elements may change is given to a
variable whose elements may change too. Elsewhere, it is refused, so that it
keeps marking only these places.

\begin{lstlisting}[style=coloredverbatimError, escapechar=@, caption=An alias to a constant]
use std::io;

fn main() {
  let dmut scores = copy [12, 15, 9];
  let seen = @\hb{alias}@ scores; // error
  println(scores, " ", seen);
}
\end{lstlisting}

The compiler says \errmsg{E4247}{aliasing the value of type mut [mut i32] to
  create a constant borrowing is prohibited}. \token{seen} cannot change the
elements, and \token{let seen = scores;} is the declaration to write. The
keyword is refused as well in front of a variable whose elements are constant,
such as one declared \token{let mut}: it has no right to change them, and
cannot give that right to another variable.

\subsection{Tuples and arrays that hold slices}

A tuple is stored in its variable, like an array, and its assignment copies its
fields. A field that is a slice is copied as a slice is, length and address, and
its elements are then shared.

\begin{lstlisting}[style=coloredverbatim]
let dmut team = (1, copy [12, 15, 9]);
let dmut other = alias team;
other.0 = 2;
other.1[0] = 20;
println(team, " ", other);
\end{lstlisting}
\vspace{-10pt}%
\begin{lstlisting}[style=bashVerb]
(1, [20, 15, 9]) (2, [20, 15, 9])
\end{lstlisting}

\token{other.0 = 2;} changes the first field of \token{other}, which is a copy:
\token{team.0} is unchanged. \token{other.1[0] = 20;} writes
into the elements that both second fields designate. \token{alias} is needed
because of that second field, whose elements both variables may change. An
array of slices behaves the same way.

\vfill\pagebreak
\section{Functions that change their arguments}
\label{sec:mutability:parameters}

A parameter holds a copy of its argument (Section~\ref{sec:functions:parameters}).
For a slice, that copy is a length and an address, and the parameter shares the
elements of the argument, as \token{seen} shared those of \token{primes}. A
function can therefore change the elements of a slice that its caller gives it,
if its parameter says so.

\lstinputlisting[style=coloredverbatim, caption={A function that changes the elements of a slice, \textit{examples/mutability/discount.yr}}, label=lst:mutability:discount]{examples/mutability/discount.yr}
\vspace{-10pt}%
\begin{lstlisting}[style=bashVerb]
[108, 72, 40]
\end{lstlisting}

The parameter \token{prices} of line~3 is declared \token{dmut}, as a variable
would be, and its type is \token{mut [mut i32]}: the function may change the
elements it receives. The call, on line~11, passes \token{alias prices}. The
parameter and the variable of \token{main} are two names for the same elements,
and \token{alias} marks the place where the second name is created, as in
Section~\ref{sec:mutability:alias}. Figure~\ref{fig:mutability:parameter} shows
the memory during the call.

\begin{figure}[h]
  \centering
  \begin{adjustbox}{max width=\linewidth}
    \begin{tikzpicture}
      \node[memcell] (mlen) at (0, 0) {3};
      \node[memaddr, right=0mm of mlen] (mptr) {};
      \node[cflab, left=3mm of mlen] (mn) {\token{prices}};
      \node[memframe, fit=(mn)(mptr)] (main) {};
      \node[memtitle] at ([xshift=2mm]main.north west) {main};

      \node[memcell] (dlen) at (0, -1.6) {3};
      \node[memaddr, right=0mm of dlen] (dptr) {};
      \node[cflab, left=3mm of dlen] (dn) {\token{prices}};
      \node[memcell] (pc) at (0, -2.6) {10};
      \node[cflab, left=3mm of pc] (pn) {\token{percent}};
      \node[memframe, fit=(dn)(dptr)(pn)] (discount) {};
      \node[memtitle] at ([xshift=2mm]discount.north west) {discount};

      \foreach \value [count=\i from 0] in {108, 72, 45} {
        \ifnum\i<2
          \node[memhit] (h\i) at (5.2 + \i * 1.0, -0.8) {\value};
        \else
          \node[memcell] (h\i) at (5.2 + \i * 1.0, -0.8) {\value};
        \fi
      }
      \draw[mempoint] (mptr.center) to[out=0, in=180] (h0.west);
      \draw[mempoint] (dptr.center) to[out=0, in=180] (h0.west);

      \node[memregion] (stack) at ($(main.north) + (0, 5mm)$) {stack};
      \node[memregion] at (h1 |- stack.base) {heap};
      \draw[memsep] (3.8, 1.0) -- (3.8, -3.3);
    \end{tikzpicture}
  \end{adjustbox}
  \caption{\label{fig:mutability:parameter} The call
    \token{discount(alias prices, 10)} of Listing~\ref{lst:mutability:discount},
    during the third turn of its loop. The call has a frame of its own, below
    that of \token{main}, and its parameter \token{prices} received the length
    and the address of the variable of \token{main}. The two designate the same
    elements: the first two are already changed.}
\end{figure}


The loop of line~4 goes through the indices of the slice, and line~5 replaces
each price with its value reduced by \token{percent} percent. The division is
between integers, so 45 reduced by 10\% gives 40, not 40.5. When the call
returns, its frame is given back, but the elements it changed are those of
\token{main}, which prints them.

The \token{alias} of the call is required. Without it, the call is refused.

\begin{lstlisting}[style=coloredverbatimError, escapechar=@, caption=A call that shares elements without \token{alias}]
use std::io;

fn discount(dmut prices: [i32], percent: i32) {
  for i in 0 .. prices.len {
    prices[i] = prices[i] * (100 - percent) / 100;
  }
}

fn main() {
  let dmut prices = copy [120, 80, 45];
  discount(@\hb{prices}@, 10); // error
  println(prices);
}
\end{lstlisting}

The compiler says \errmsg{E4226}{the call operator is not defined for
  \errhole{} and \errhole{}}, since no function \token{discount} accepts the
arguments as they are written, and the note under it gives the reason,
\textit{implicit alias of type mut [mut i32] is prohibited}. A reader of
\token{main} sees from the call alone that \token{discount} may change
\token{prices}: the \token{alias} is there. A call without it, such as
\token{average(prices)}, gives the function no right to change the elements.

\subsection{The parameter itself is still a constant}

\token{dmut} lets the function change the elements, not the parameter: giving
\token{prices} a new slice inside \token{discount}, \token{prices = copy [0];},
is refused with \errmsg{E4156}{prices is a value parameter, and cannot be used
  as a lvalue}. The new slice would only replace the copy of the length and the
address held by the frame of \token{discount}, and would have no effect on the
caller. Section~\ref{sec:mutability:references} shows how a function gives a new
value to a variable of its caller.

For the same reason, \token{mut} and \token{dmut} are refused on a parameter
whose value borrows nothing, such as \token{dmut n: i32} or
\token{dmut days: [i32 ; 12]}. Its argument is copied whole into the frame of
the function, so changing it could never affect the caller, and the compiler
says \errmsg{E4138}{a parameter cannot be mutable, if it is not a reference or
  does not borrow mutable data}.

\subsection{Returning elements that can change}

A function that creates a slice can return it with elements that can change,
with a return type written \token{dmut [i32]}.

\begin{lstlisting}[style=coloredverbatim]
fn zeros(n: usize)-> dmut [i32] {
  copy [0 ; n]
}

fn main() {
  let dmut counts = zeros(3);
  counts[1] = 5;
  println(counts);
}
\end{lstlisting}
\vspace{-10pt}%
\begin{lstlisting}[style=bashVerb]
[0, 5, 0]
\end{lstlisting}

The result of \token{copy} shares its elements with nothing, so it needs no
\token{alias}. Neither does the call, on the side of \token{main}: the function
has returned, and only \token{counts} designates the elements. A function that
returns the elements of one of its variables declared \token{dmut}, however,
returns them with \token{alias}, as in \token{alias result}: the variable and
the value returned designate the same elements, however briefly.

\vfill\pagebreak
\section{Copies: \texttt{copy} and \texttt{dcopy}}
\label{sec:mutability:copies}

Sharing is the default, and a program that wants elements of its own asks for a
copy. \token{copy} was met as the way to make a slice from an array
(Section~\ref{sec:collections:copy}). It applies to a slice as well: it
reserves a new block, copies the elements into it, and gives a slice that
designates the new block. The next listing keeps the prices of
Listing~\ref{lst:mutability:discount} before the discount changes them.

\lstinputlisting[style=coloredverbatim, caption={Keeping the elements before a change, \textit{examples/mutability/backup.yr}}, label=lst:mutability:backup]{examples/mutability/backup.yr}
\vspace{-10pt}%
\begin{lstlisting}[style=bashVerb]
Now:    [108, 72, 40]
Seen:   [108, 72, 40]
Before: [120, 80, 45]
\end{lstlisting}

\token{seen}, on line~11, shares the elements of \token{prices}, and sees the
discount. \token{before}, on line~12, has elements of its own, and keeps the
prices as they were. A copy shares nothing with the original, so it needs no
\token{alias}, even to a variable declared \token{dmut}:
\token{let dmut other = copy prices;} is accepted, and changing \token{other}
leaves \token{prices} as it is.

A copy has a cost that sharing does not have: an allocation, and the time to
copy every element. A slice of a million elements passed to a function costs
two numbers, and its copy a million elements. A program copies when it needs
elements of its own, not by default.

\subsection{Copying a slice of slices}

The elements of a slice of slices are slices, and \token{copy} copies them as
slices are copied: a length and an address each. The copy has a new block of
rows, but its rows designate the elements of the rows of the original.
\token{dcopy}, for \textit{deep copy}, copies every level: the block of rows,
and each row with its elements.

\begin{lstlisting}[style=coloredverbatim]
let dmut grid = dcopy [[1, 2], [3, 4]];
let shallow = copy grid;
let dmut deep = dcopy grid;
grid[0][0] = 9;
deep[1][1] = 0;
println(grid, " ", shallow, " ", deep);
\end{lstlisting}
\vspace{-10pt}%
\begin{lstlisting}[style=bashVerb]
[[9, 2], [3, 4]] [[9, 2], [3, 4]] [[1, 2], [3, 0]]
\end{lstlisting}

\token{grid[0][0] = 9;} writes into the first row of \token{grid}, which is also
the first row of \token{shallow}. \token{deep} has rows of its own, and neither
sees the change of the other. Figure~\ref{fig:mutability:copies} shows the three
variables. The first line uses \token{dcopy} as well: \token{[[1, 2], [3, 4]]}
is an array of arrays, and \token{dcopy} turns each level into a slice, the
rows included.

\begin{figure}[h]
  \centering
  \begin{adjustbox}{max width=\linewidth}
    \begin{tikzpicture}
      \foreach \name/\y in {grid/0, shallow/-2.4, deep/-4.0} {
        \node[memcell] (\name len) at (0, \y) {2};
        \node[memaddr, right=0mm of \name len] (\name ptr) {};
        \node[memcell] (\name r0len) at (4.8, \y) {2};
        \node[memaddr, right=0mm of \name r0len] (\name r0ptr) {};
        \node[memcell, right=0mm of \name r0ptr] (\name r1len) {2};
        \node[memaddr, right=0mm of \name r1len] (\name r1ptr) {};
        \draw[mempoint] (\name ptr.center) -- (\name r0len.west);
      }
      \node[cflab, left=3mm of gridlen] (gridn) {\token{grid}};
      \node[cflab, left=3mm of shallowlen] (shallown) {\token{shallow}};
      \node[cflab, left=3mm of deeplen] (deepn) {\token{deep}};
      \node[memframe, fit=(gridn)(shallown)(gridptr)(deepptr)] (main) {};
      \node[memtitle] at ([xshift=2mm]main.north west) {main};

      \node[memcell] (a0) at (5.8, -1.2) {1};
      \node[memcell, right=0mm of a0] (a1) {2};
      \node[memcell] (b0) at (8.3, -1.2) {3};
      \node[memcell, right=0mm of b0] (b1) {4};
      \node[memcell] (c0) at (5.8, -5.2) {1};
      \node[memcell, right=0mm of c0] (c1) {2};
      \node[memcell] (e0) at (8.3, -5.2) {3};
      \node[memcell, right=0mm of e0] (e1) {4};

      \draw[mempoint] (gridr0ptr.center) to[out=-90, in=90] (a0.north);
      \draw[mempoint] (gridr1ptr.center) to[out=-90, in=90] (b0.north);
      \draw[mempoint] (shallowr0ptr.center) to[out=90, in=-90] (a0.south);
      \draw[mempoint] (shallowr1ptr.center) to[out=90, in=-90] (b0.south);
      \draw[mempoint] (deepr0ptr.center) to[out=-90, in=90] (c0.north);
      \draw[mempoint] (deepr1ptr.center) to[out=-90, in=90] (e0.north);

      \node[memregion] (stack) at ($(main.north) + (0, 5mm)$) {stack};
      \node[memregion] at ({$(gridr0len)!0.5!(b1)$} |- stack.base) {heap};
      \draw[memsep] (3.3, 1.2) -- (3.3, -5.8);
    \end{tikzpicture}
  \end{adjustbox}
  \caption{\label{fig:mutability:copies} A slice of two rows, \token{grid}, and
    two copies of it. \token{shallow = copy grid} has a block of its own for
    the two rows, but each row is copied as a slice is, length and address,
    and still designates the elements of \token{grid}. \token{deep = dcopy grid}
    copies the rows as well: it shares nothing with \token{grid}.}
\end{figure}


Since the rows of \token{shallow} are those of \token{grid}, whose elements can
change, \token{shallow} cannot be declared \token{dmut}. The type of
\token{copy grid} is \token{mut [mut [i32]]}: the rows are constant through the
copy, and \token{let dmut shallow = copy grid;} is refused with
\errmsg{E4045}{discarding the constant property is prohibited}, the note under
it saying \textit{from mut [mut [i32]] to mut [mut [mut i32]]}. A copy made
with \token{copy} can therefore never change the rows of the original. To
change the rows of a copy, the copy is made with \token{dcopy}.

\subsection{How deep \texttt{dcopy} goes}

\token{dcopy} is defined in terms of itself. The deep copy of a slice is a new
block whose elements are the deep copies of the elements of the slice. When the
elements are slices in turn, each is deep-copied the same way, and so on, until
the elements borrow nothing: an integer is its own deep copy. Written by hand
for a grid, one level of that recursion is a \token{copy} of each row.

\begin{lstlisting}[style=coloredverbatim]
fn deepCopy(grid: [[i32]])-> dmut [[i32]] {
  let dmut result: [[i32]] = [];
  for row in grid {
    result ~= [copy row];
  }
  alias result
}
\end{lstlisting}

\token{dcopy} does the same at every level, however many there are. A cube of
numbers, a slice of grids of type \token{[[[i32]]]}, has three levels, and
\token{dcopy cube} copies the block of grids, the block of rows of each grid,
and the elements of each row.

\begin{lstlisting}[style=coloredverbatim]
let dmut cube = dcopy [[[1, 2], [3, 4]], [[5, 6], [7, 8]]];
let dmut other = dcopy cube;
other[1][1][1] = 0;
println(cube, " ", other);
\end{lstlisting}
\vspace{-10pt}%
\begin{lstlisting}[style=bashVerb]
[[[1, 2], [3, 4]], [[5, 6], [7, 8]]] [[[1, 2], [3, 4]], [[5, 6], [7, 0]]]
\end{lstlisting}

\subsubsection*{A value that contains itself}

A recursion needs a base case (Section~\ref{sec:functions:infinite_recursion}),
and that of \token{dcopy} is the element that borrows nothing. A slice that
contained itself, directly or through another slice, would never reach it: the
deep copy of its elements would include its own deep copy, which would include
it again, and so on until the stack overflows. The types prevent such a slice
from existing.

\begin{lstlisting}[style=coloredverbatimError, escapechar=@, caption=A slice that would contain itself]
use std::io;

fn main() {
  let dmut b = copy [copy [1]];
  let dmut a = copy [alias b];
  b[0] = @\hb{a}@; // error
  println(a);
}
\end{lstlisting}

The compiler says \errmsg{E4095}{incompatible types mut [mut i32] and mut [mut
  [mut [mut i32]]]}. \token{b} is a slice of slices of integers, and \token{a}
a slice of such slices: one level more. For \token{b[0] = a;} to be accepted,
an element of \token{b} would have to have the type of \token{a}, which holds
the type of \token{b}, which holds that of its elements, without end. The type
of a slice is always one level deeper than that of its elements, so no slice,
array or map can hold itself, and \token{dcopy} always reaches the bottom.

The classes of Chapter~\ref{chap:custom_types} can refer to each other, and to
themselves through others: two objects each naming the other form a cycle. The
compiler therefore does not deep-copy a class by itself: \token{dcopy} on an
object whose class does not say how it is copied is refused with
\errmsg{E4018}{class type \errhole{} does not implement trait
  core::types::Copiable}, and the class decides what its deep copy does with a
cycle.

\subsubsection*{Maps and tuples}

A map is deep-copied the same way, its values included: the deep copy of a map
from names to slices of grades has slices of its own. Tuples are the exception.
Neither \token{copy} nor \token{dcopy} applies to a tuple, which is refused with
\errmsg{E4167}{no copy exists for type (i32, [i32])}, and therefore not to a
slice whose elements are tuples holding slices either. Such a value is copied
field by field.

\subsection{Who makes the copy}

A function that changes a slice can be written in two ways. The first changes
the elements it receives, in place, and the second copies them first and
returns the copy, changed. The next listing sorts a slice in increasing order,
in place: it looks for the smallest element and exchanges it with the first,
then for the smallest of the others, exchanged with the second, and so on. This
method is called the \textit{selection sort}.

\lstinputlisting[style=coloredverbatim, caption={Sorting in place, \textit{examples/mutability/sort.yr}}, label=lst:mutability:sort]{examples/mutability/sort.yr}
\vspace{-10pt}%
\begin{lstlisting}[style=bashVerb]
[3, 7, 19, 25, 42]
[42, 7, 19, 3, 25] [3, 7, 19, 25, 42]
\end{lstlisting}

\token{sort} takes the slice as a \token{dmut} parameter, and makes no copy.
Its caller chooses. On line~23, \token{main} sorts \token{scores} itself, and
pays for no copy. On lines~27 and~28, it wants to keep \token{arrivals} in the
order of arrival, so it makes the copy, and sorts the copy.

The other way puts the copy inside the function.

%% check: skip
\begin{lstlisting}[style=coloredverbatim]
fn sorted(values: [i32])-> [i32] {
  let dmut result = copy values;
  sort(alias result);
  result
}
\end{lstlisting}

\token{sorted} never changes its argument, but every call copies it, including
\token{scores = sorted(scores);}, which throws the old elements away at once.
The caller of \token{sorted} cannot avoid the copy, whereas the caller of
\token{sort} makes one only when it needs one. A function that changes
elements is therefore best written in place, with a \token{dmut} parameter, and
the decision to copy left to the caller, who knows whether the original is
still needed.

The copy belongs inside the function when its contract says that it only reads:
a function that computes the median of a slice, the element in the middle once
the slice is sorted, takes a plain parameter, and must not sort the slice of its
caller. It sorts a copy, as Exercise~8 does.

\subsection{The repeated row}
\label{sec:mutability:repeated_row}

A caution of Section~\ref{sec:collections:grids} warned against
\token{copy [copy [0 ; cols] ; rows]}, which looks like a grid of \token{rows}
rows but changes all of them at once. The reason is now visible. \token{[v ; n]}
computes \token{v} once, and makes \token{n} copies of the value: for a slice,
\token{n} copies of the same length and the same address.
Figure~\ref{fig:mutability:repeated_row} shows the grid.

\begin{figure}[h]
  \centering
  \begin{adjustbox}{max width=\linewidth}
    \begin{tikzpicture}
      \node[memcell] (glen) at (0, 0) {3};
      \node[memaddr, right=0mm of glen] (gptr) {};
      \node[cflab, left=3mm of glen] (gn) {\token{grid}};
      \node[memframe, fit=(gn)(gptr)] (main) {};
      \node[memtitle] at ([xshift=2mm]main.north west) {main};

      \node[memcell] (r0len) at (4.8, 0) {4};
      \node[memaddr, right=0mm of r0len] (r0ptr) {};
      \node[memcell, right=0mm of r0ptr] (r1len) {4};
      \node[memaddr, right=0mm of r1len] (r1ptr) {};
      \node[memcell, right=0mm of r1ptr] (r2len) {4};
      \node[memaddr, right=0mm of r2len] (r2ptr) {};
      \draw[mempoint] (gptr.center) -- (r0len.west);
      \foreach \i in {0, 1, 2} {
        \node[memnote, above=0.5mm of r\i len.north east] {row \i};
      }

      \foreach \value [count=\i from 0] in {0, 0, 7, 0} {
        \ifnum\i=2
          \node[memhit] (e\i) at (7.2 + \i * 1.0, -1.9) {\value};
        \else
          \node[memcell] (e\i) at (7.2 + \i * 1.0, -1.9) {\value};
        \fi
      }
      \foreach \i in {0, 1, 2} {
        \draw[mempoint] (r\i ptr.center) to[out=-90, in=90] (e0.north);
      }

      \node[memregion] (stack) at ($(main.north) + (0, 5mm)$) {stack};
      \node[memregion] at ({$(r0len)!0.5!(r2ptr)$} |- stack.base) {heap};
      \draw[memsep] (3.3, 0.9) -- (3.3, -2.6);
    \end{tikzpicture}
  \end{adjustbox}
  \caption{\label{fig:mutability:repeated_row} The grid
    \token{copy [copy [0 ; 4] ; 3]}. \token{copy [0 ; 4]} made one row, and
    \token{[v ; 3]} repeated that value three times: three copies of the same
    length and the same address. \token{grid[1][2] = 7;} changes the element
    that the three rows designate.}
\end{figure}


\begin{lstlisting}[style=coloredverbatim]
let dmut grid = copy [copy [0 ; 4] ; 3];
grid[1][2] = 7;
println(grid);
\end{lstlisting}
\vspace{-10pt}%
\begin{lstlisting}[style=bashVerb]
[[0, 0, 7, 0], [0, 0, 7, 0], [0, 0, 7, 0]]
\end{lstlisting}

The outer \token{copy} copies the three rows as slices, and does not separate
them. The loop of Section~\ref{sec:collections:grids} makes a new row at each
turn, with a \token{copy} of its own, and gives distinct rows. \token{dcopy} of
the grid above also gives distinct rows, since it copies each of them in turn:
\token{let dmut fixed = dcopy grid;} is a grid of three rows of its own, which
start with the 7 of the shared row. In version \gycversion{} of the compiler,
\token{dcopy} written directly around the repeated literal,
\token{dcopy [copy [0 ; 4] ; 3]}, still gives a single shared row, and the loop
is the way to write it.

\vfill\pagebreak
\section{References}
\label{sec:mutability:references}

\token{discount} changes the elements of its caller's slice, because the
elements are shared. It cannot change a variable of its caller: the parameter
is a copy of the argument, and nothing it does reaches the variable the argument
came from. A function that exchanges the values of two integers needs more: it
must write into the variables of its caller.

A \textit{reference} is a parameter that holds the address of a variable of the
caller rather than a copy of its value. Reading the parameter reads that
variable, and assigning the parameter writes into it. It is declared with
\token{ref} before the name, and \token{ref mut} when the function may write.

\lstinputlisting[style=coloredverbatim, caption={Exchanging two variables, \textit{examples/mutability/swap.yr}}, label=lst:mutability:swap]{examples/mutability/swap.yr}
\vspace{-10pt}%
\begin{lstlisting}[style=bashVerb]
x = 2, y = 1
\end{lstlisting}

The parameters \token{a} and \token{b} of line~3 are references to integers
that the function may change. In the body, they are used as ordinary variables:
line~4 reads the value of the variable that \token{a} refers to, and line~5
writes into it. The call, on line~12, passes \token{ref x} and \token{ref y}:
the addresses of the two variables of \token{main}, which must be declared
\token{mut} for the function to change them.
Figure~\ref{fig:mutability:ref_parameter} shows the call.

\begin{figure}[h]
  \centering
  \begin{adjustbox}{max width=\linewidth}
    \begin{tikzpicture}
      \node[memhit] (x) at (0, 0) {2};
      \node[cflab, left=3mm of x] (xn) {\token{x}};
      \node[memcell] (y) at (0, -1.0) {2};
      \node[cflab, left=3mm of y] {\token{y}};
      \node[memframe, fit=(xn)(x)(y)] (main) {};
      \node[memtitle] at ([xshift=2mm]main.north west) {main};

      \node[memaddr] (a) at (0, -2.6) {};
      \node[cflab, left=3mm of a] (an) {\token{a}};
      \node[memaddr] (b) at (0, -3.6) {};
      \node[cflab, left=3mm of b] {\token{b}};
      \node[memcell] (old) at (0, -4.6) {1};
      \node[cflab, left=3mm of old] (on) {\token{old}};
      \node[memframe, fit=(an)(a)(old)] (swap) {};
      \node[memtitle] at ([xshift=2mm]swap.north west) {swap};

      \draw[mempoint] (a.center) to[out=0, in=0, looseness=1.2] (x.east);
      \draw[mempoint] (b.center) to[out=0, in=0, looseness=1.2] (y.east);

      \node[memregion] at ($(main.north) + (0, 5mm)$) {stack};
    \end{tikzpicture}
  \end{adjustbox}
  \caption{\label{fig:mutability:ref_parameter} The call
    \token{swap(ref x, ref y)} of Listing~\ref{lst:mutability:swap}, after
    \token{a = b;}. The parameters \token{a} and \token{b} hold the addresses of
    the variables of \token{main}, and \token{a = b;} wrote into \token{x}.
    \token{old} still holds the first value of \token{x}, which
    \token{b = old;} writes into \token{y} next.}
\end{figure}


As \token{alias} does, the \token{ref} of the call tells a reader of
\token{main} that the call may change the variables it names. It is required.
\token{swap(x, y)} is refused with \errmsg{E4226}{the call operator is not
  defined for \errhole{} and \errhole{}}, the note under it saying
\textit{cannot pass value of type mut i32 as ref}. A reference refers to a
variable, so its argument must be one: \token{swap(ref x, ref 3)} is refused
with \errmsg{E4151}{not a lvalue}, and a variable declared without \token{mut}
cannot be passed to \token{ref mut}, since the function would change a
constant.

\subsection{Giving the caller's variable a new slice}

A reference can refer to a variable of any type, including a slice. Through
it, a function can give its caller's variable a new slice, which a
\token{dmut} parameter cannot do (Section~\ref{sec:mutability:parameters}).

\lstinputlisting[style=coloredverbatim, caption={Appending to the caller's slice, \textit{examples/mutability/history.yr}}, label=lst:mutability:history]{examples/mutability/history.yr}
\vspace{-10pt}%
\begin{lstlisting}[style=bashVerb]
[18, 21, 19]
\end{lstlisting}

\tokennolst{history \mtildeascii{}= [value]}, on line~4, gives the variable a
new slice, one element longer. With a parameter declared \token{mut history:
  [i32]}, this would have changed a copy, and \token{main} would have printed an
empty slice; the compiler refuses such a parameter anyway, with the
\errmsg{E4138}{a parameter cannot be mutable, if it is not a reference or does
  not borrow mutable data} of Section~\ref{sec:mutability:parameters}. The
variable of \token{main} is \token{mut}: only the variable changes, not the
elements already in it.

\subsection{Constant references}

A reference without \token{mut} gives read access to the caller's variable. It
is useful for a large value stored in its variable, such as an array of a
thousand elements: a parameter of type \token{[i32 ; 1000]} copies the thousand
elements into the frame at each call, a reference copies an address.

\begin{lstlisting}[style=coloredverbatim]
fn last(ref values: [i32 ; 1000])-> i32 {
  values[999]
}

fn main() {
  let big = [7 ; 1000];
  println(last(big));
}
\end{lstlisting}

The call writes nothing: the function cannot change the variable, so there is
nothing for the reader to be warned of. The \token{ref} keyword in a call is
written only for a reference that may change, and is refused elsewhere, as
\token{alias} is. A slice needs no constant reference: it is already two
numbers, and a plain parameter shares its elements.

\subsection{References to elements}

The argument of a reference can be an element of an array or of a slice, which
is a place in memory as a variable is. \token{swap(ref t[0], ref t[2])}
exchanges the first and the third elements of \token{t}, a slice declared
\token{dmut}. To change the elements of an array of its caller, a function
takes a reference to the whole array, declared \token{ref dmut}: \token{ref mut}
alone lets it give the array a new value, but not change its elements, for the
reason of Section~\ref{sec:mutability:levels}.

\begin{lstlisting}[style=coloredverbatim]
fn clear(ref dmut counts: [i32 ; 6]) {
  for i in 0 .. 6 {
    counts[i] = 0;
  }
}

fn main() {
  let dmut counts = [1, 2, 3, 4, 5, 6];
  clear(ref counts);
  println(counts);
}
\end{lstlisting}

\subsection{Choosing a parameter}

\begin{center}
  \begin{adjustbox}{max width=\linewidth}
    \begin{tabular}{|l|l|l|}
      \hline
      The function & Parameter & Call\\[0pt]
      \hline
      \hline
      reads its argument & \token{x: T} & \token{f(x)}\\[0pt]
      reads a large array without copying it & \token{ref x: [T ; n]} & \token{f(x)}\\[0pt]
      changes the elements of a slice or a map & \token{dmut x: [T]} & \token{f(alias x)}\\[0pt]
      gives the caller's variable a new value & \token{ref mut x: T} & \token{f(ref x)}\\[0pt]
      changes the elements of the caller's array & \token{ref dmut x: [T ; n]} & \token{f(ref x)}\\[0pt]
      \hline
    \end{tabular}
  \end{adjustbox}
\end{center}

A reference is only a parameter, or a local variable declared
\token{let ref}, which Part II describes (Section~\ref{sec:memory:ref_variable}).
It cannot be stored in another value, nor returned by a function: the variable
it refers to may be in the frame of the function, which disappears when the
function returns.

\vfill\pagebreak
\section{Loops that change the elements}
\label{sec:mutability:loops}

The iterator of a \token{for} loop holds a copy of the current element, and is a
constant (Section~\ref{sec:collections:arrays}). To change the elements while
going through them, Chapter~\ref{chap:collections} went through the indices and
wrote \token{counts[i]}. An iterator can instead be a reference to the current
element, as a parameter can be a reference to a variable.

\lstinputlisting[style=coloredverbatim, caption={Raising every grade in place, \textit{examples/mutability/raise.yr}}, label=lst:mutability:raise]{examples/mutability/raise.yr}
\vspace{-10pt}%
\begin{lstlisting}[style=bashVerb]
[14, 20, 10, 17]
\end{lstlisting}

On line~5, \token{ref mut grade} declares the iterator as a reference that may
change, and \token{ref grades} gives the loop the right to change the elements
of \token{grades}, as \token{ref x} gave it to \token{swap}. At each turn,
\token{grade} refers to one element of the slice, and line~6 writes into that
element: two points more, but no more than 20. With an index, the reference is
the second iterator, \token{for i, ref mut x in ref values}. The same loop works
on an array declared \token{dmut}.

The two keywords go together, and each is refused without the other. Without
\token{ref}, the iterator is a copy, and a copy that could change would change
nothing: \token{for mut grade in grades} is refused with
\errmsg{E4137}{an iterator cannot be mutable, if it is not a reference or does
  not borrow mutable data}. An assignment to a constant iterator, as in
\token{for grade in grades \{ grade = grade + 2; \}}, is refused with
\errmsg{E4153}{grade is a value iterator, and cannot be used as a lvalue}.

\token{ref grades} refers to a variable, and a parameter does not count as one
here: over a parameter declared \token{dmut values: [i32]}, the loop
\token{for ref mut v in ref values} is refused with \errmsg{E4156}{values is a
  value parameter, and cannot be used as a lvalue}, although only the elements
would change. A function that changes the elements of its parameter goes
through the indices, as \token{discount} does in
Listing~\ref{lst:mutability:discount}.

\subsection{Going through the rows of a grid}

The rows of a slice of slices are slices, and a loop over the grid gives a
copy of each row, a length and an address. That copy shares the elements of the
row, as a parameter does, so a loop can change the elements of the rows without
a reference: the iterator is declared \token{dmut}, and the grid is given with
\token{alias}, as a slice is given to a \token{dmut} parameter.

\begin{lstlisting}[style=coloredverbatim]
let dmut grid = dcopy [[1, 2], [3, 4]];
for dmut row in alias grid {
  for j in 0 .. row.len {
    row[j] = row[j] * 10;
  }
}
println(grid);
\end{lstlisting}
\vspace{-10pt}%
\begin{lstlisting}[style=bashVerb]
[[10, 20], [30, 40]]
\end{lstlisting}

\token{row} is still a copy: \token{row = copy [0, 0];} would give a new slice
to the iterator, not to the grid, and is refused. To replace whole rows, the
iterator is a reference, \token{for ref mut row in ref grid}.

The iterator of the values of a map cannot be a reference, nor mutable. To
change the values of a map in a loop, the loop writes through the key, as
Listing~\ref{lst:collections:letters} writes \token{counts[c]}.

\begin{lstlisting}[style=coloredverbatim]
let dmut stock = copy ["apples" => 3, "pears" => 5];
for fruit, count in stock {
  stock[fruit] = count * 2;
}
println(stock);
\end{lstlisting}
\vspace{-10pt}%
\begin{lstlisting}[style=bashVerb]
[apples => 6, pears => 10]
\end{lstlisting}

\vfill\pagebreak
\section{Patterns that change what they match}
\label{sec:mutability:patterns}

A pattern declares variables: the names of \token{(d, m)} or of
\token{Ok(n)} (Section~\ref{sec:collections:match}). Like the iterators of a
loop, they are constant copies of the parts of the value matched. They can also
be declared as a loop iterator can, with \token{dmut} for a part whose elements
change, and with \token{ref mut} for a reference to the part itself. The value
matched then carries \token{alias} or \token{ref}, as the argument of a call
does.

\subsection{\texttt{match alias}}

A variable of a pattern that shares elements that it may change is declared
\token{dmut}, and the value matched is given with \token{alias}.

\begin{lstlisting}[style=coloredverbatim]
let dmut team = (1, copy [12, 15, 9]);
match alias team {
  (1, dmut scores) => {
    scores[0] = 20;
  }
  _ => {}
}
println(team);
\end{lstlisting}
\vspace{-10pt}%
\begin{lstlisting}[style=bashVerb]
(1, [20, 15, 9])
\end{lstlisting}

\token{scores} is a copy of the second field of \token{team}, a length and an
address, and shares its elements as the tuples of
Section~\ref{sec:mutability:alias} did. Without \token{alias}, the pattern is
refused with \errmsg{E4045}{discarding the constant property is prohibited},
and the note under it names the field, \textit{for pattern field access 1 on
  type (mut i32, mut [mut i32])}.

\subsection{\texttt{match ref}}

A variable of a pattern declared \token{ref mut} is a reference to a part of
the value matched, and assigning it writes into that part. The value matched is
given with \token{ref}, and must be a variable, or a part of one. The next
listing puts the two bounds of an interval in order.

\begin{lstlisting}[style=coloredverbatim]
let dmut range = (9, 4);
match ref range {
  (ref mut low, ref mut high) if low > high => {
    let old = low;
    low = high;
    high = old;
  }
  _ => {}
}
println(range);
\end{lstlisting}
\vspace{-10pt}%
\begin{lstlisting}[style=bashVerb]
(4, 9)
\end{lstlisting}

\token{low} and \token{high} refer to the two fields of \token{range}, and the
guard reads them as it would read copies. \token{range} is declared
\token{dmut}: its fields are a level of their own, as the elements of an array
are, and \token{let mut} would let the program give \token{range} a new tuple,
but not change one field.

The same works on an option, whose value can then be changed in place, and on
a slice, whose pattern can take references to its first elements.

\begin{lstlisting}[style=coloredverbatim]
let dmut best: i32? = none;
for score in copy [12, 17, 9] {
  if let Ok(ref mut b) = ref best {
    if score > b {
      b = score;
    }
  } else {
    best = score?;
  }
}
println(best);

let dmut word = copy "hello";
match ref word {
  [ref mut first, _...] => {
    first = 'H';
  }
  _ => {}
}
println(word);
\end{lstlisting}
\vspace{-10pt}%
\begin{lstlisting}[style=bashVerb]
Ok(17)
Hello
\end{lstlisting}

A \token{let} declaration is a pattern too, and takes the same keywords. The
next declaration makes two references to the fields of \token{range}.

\begin{lstlisting}[style=coloredverbatim]
let dmut range = (9, 4);
let (ref mut low, ref mut high) = ref range;
low = 0;
println(range, " ", high);
\end{lstlisting}
\vspace{-10pt}%
\begin{lstlisting}[style=bashVerb]
(0, 4) 4
\end{lstlisting}

\subsection{Which keyword}

A pattern can mix the three kinds of variables: constant copies, \token{dmut}
copies that share elements, and \token{ref mut} references. The keyword in front
of the value matched depends on the variables of all the arms.

\begin{itemize}
\item \token{ref}, if a variable is declared \token{ref mut}. It also covers the
  \token{dmut} variables of the same pattern, as in
  \token{[ref mut first, dmut rest...]}, which takes a reference to the first
  element and shares the others.
\item \token{alias}, if a variable is declared \token{dmut} and shares elements
  that it may change, and none is a reference that may change.
\item Neither, if every variable is a constant. \token{match ref range} with the
  pattern \token{(a, b)} is refused with \errmsg{E4252}{referencing the value
    has no effect}, the note under it saying \textit{there is no pattern that
    requires the value to be passed by mutable reference}, and
  \token{match alias} is refused in the same case with the
  \errmsg{E4247}{aliasing the value \errhole{} to create a constant borrowing
    is prohibited} of Section~\ref{sec:mutability:alias}.
\end{itemize}

A variable declared \token{ref mut} when the value is matched without
\token{ref} is refused with \errmsg{E4089}{ref of type mut i32 must be mutable,
  but is const}: without the keyword, the parts of the value are constant for
the pattern.

\vfill\pagebreak
\section{A game of life}
\label{sec:mutability:life}

The \textit{game of life}, invented by the mathematician John Conway, is played
by nobody: a grid of cells, each alive or dead, changes by itself from one
generation to the next, according to two rules that count the living
neighbours of each cell among the eight around it.

\begin{itemize}
\item A living cell stays alive with two or three living neighbours, and dies
  otherwise, of loneliness or of overcrowding.
\item A dead cell comes to life with exactly three living neighbours.
\end{itemize}

From such simple rules come shapes that last, that repeat, or that travel. The
next program follows a \textit{glider}, five cells that move one cell
diagonally every four generations. The board is changed in place: a single
grid, which each generation rewrites.

\lstinputlisting[style=coloredverbatim, caption={The game of life, \textit{examples/mutability/life.yr}}, label=lst:mutability:life]{examples/mutability/life.yr}
\vspace{-10pt}%
\begin{lstlisting}[style=bashVerb, escapechar=@+]
Generation 0
.#....
..#...
###...
......
......
Generation 1
......
#.#...
.##...
.#....
......
@+\textit{(12 more lines)}@+
Generation 4
......
..#...
...#..
.###..
......
\end{lstlisting}

After four generations, the glider has the shape it started with, one row lower
and one column further right.

\subsection{Building the board}

A board is a slice of rows, and a row a slice of booleans, \token{true} for a
living cell: a \token{[[bool]]}. \token{parse} builds it from a picture, one
string per row, where \tokennolst{'\#'} stands for a living cell. \token{main} gives
it the picture as a slice, hence the \token{copy} of line~80.

Each row is built by appending to \token{row}, declared \token{dmut} on
line~13 since the program will change its cells. Line~17 appends the row to the
board, and the row becomes an element of \token{board}. The variable \token{row}
and the element then designate the same cells, both of which may change them,
hence \token{alias row}. The function returns the board with \token{alias}, on
line~19, for the reason of Section~\ref{sec:mutability:parameters}, and
\token{main} receives it in a variable declared \token{dmut}.

\subsection{Counting the neighbours}

\token{neighbours} looks at the three rows around the cell, \token{i - 1},
\token{i} and \token{i + 1}, and the three columns around it, but skips the
cell itself, on line~40. The rows are \token{usize}, which cannot be negative,
so the row before row 0 cannot be computed as \token{i - 1}. The board wraps
around instead: the row before the first is the last, and the row after the
last is the first. Line~38 adds \token{rows} before going back one row, and
takes the remainder of the division by \token{rows}, which brings the result
back between 0 and \token{rows - 1}. The iterators \token{di} and \token{dj} go
from 0 to 2, and are \token{usize} as well, with the suffix \token{us}, since
Ymir does not mix integer types in an operation.

\subsection{The next generation, in place}

\token{step} is the heart of the program. It receives the board as a
\token{dmut} parameter, called with \token{alias board} on line~88, and rewrites
its cells: \token{main} sees the next generation when the call returns. The
loop of line~56 gives each row as an iterator declared \token{dmut}, as in
Section~\ref{sec:mutability:loops}, and line~59 writes the new state of each
cell.

The state of a cell depends on its neighbours in the \textit{previous}
generation. Counting them on the board itself would go wrong: by the time the
loop reaches a cell, the cells before it already hold the new generation. Line~55
keeps the previous generation in \token{before}, and the loop counts on it.
The copy must be a \token{dcopy}. With \token{copy}, \token{before} would have
a block of rows of its own, but its rows would still designate the cells of the
board (Section~\ref{sec:mutability:copies}), and would change with them.
Exercise~6 tries it.

\token{neighbours} and \token{show} take the board as a plain parameter: they
only read it. \token{parse} returns a board whose cells can change, and
\token{step} changes one, and their signatures say so. So does \token{main}:
its only \token{alias} is the call of line~88, the one place where the board
changes.

\vfill%
\pagebreak

\section*{Exercises}
\markright{\MakeUppercase{Exercises}}%
\subsection*{Exercise 1}

What does the following program print? Work it out by hand, drawing the memory
as the figures of this chapter do, then check your answer by compiling and
running it.

\begin{lstlisting}[style=coloredverbatim]
use std::io;

fn main() {
  let dmut a = copy [1, 2, 3];
  let b = a;
  let dmut c = alias a;
  let dmut d = copy a;
  c[0] = 10;
  d[1] = 20;
  a = copy [7, 8, 9];
  a[2] = 30;
  println(a, " ", b, " ", c, " ", d);
}
\end{lstlisting}

\subsection*{Exercise 2}

The following program should set the three counts to 0, and multiply the two
prices by 3. It contains three errors. Find them, correct them, and compile the
program to test your corrections.

\begin{lstlisting}[style=coloredverbatimError, caption=Ymir program with errors]
use std::io;

fn reset(dmut values: [i32 ; 3]) {
  for i in 0 .. 3 {
    values[i] = 0;
  }
}

fn scale(dmut values: [i32], factor: i32) {
  for mut v in values {
    v = v * factor;
  }
}

fn main() {
  let dmut counts = [4, 5, 6];
  reset(counts);
  let dmut prices = copy [10, 20];
  scale(prices, 3);
  println(counts, " ", prices);
}
\end{lstlisting}

\subsection*{Exercise 3}

Exercise~1 of Chapter~\ref{chap:collections} returned a new slice holding the
elements of another in the reverse order. Write instead a function
\token{reverseInPlace(dmut values: [i32])} that reverses the elements of the
slice it receives, without creating a new one, using the function \token{swap}
of Listing~\ref{lst:mutability:swap}. Check it on a slice of five elements, one
of four, and the empty slice.

\subsection*{Exercise 4}

Write a function \token{removeAll(ref mut values: [i32], unwanted: i32)} that
removes every occurrence of \token{unwanted} from the slice of its caller: after
the call, the variable passed holds only the other elements, in the same order.
Keep a copy of the variable made with \token{let before = readings;} before the
call, and print it after. Does it change? Why?

\subsection*{Exercise 5}

A shop keeps its stock in a map from the name of each item to the number of
items left. Write a function \token{deliver} that adds a delivery, a map of the
same type, to the stock of its caller. An item that the stock does not know yet
is added to it. Test it with a stock of 12 apples and 4 pears, and a delivery of
6 pears and 20 plums.

\subsection*{Exercise 6}

In Listing~\ref{lst:mutability:life}, replace the \token{dcopy} of line~55 with
\token{copy}, then with nothing at all, \token{let before = board;}. Both
programs compile. What do they print, and why?

\subsection*{Exercise 7}

Some shapes of the game of life stop changing after a few generations. Change
\token{step} so that it returns whether the board changed, and write a
\token{main} that steps a row of four living cells on a board of 7 by 7 until it
stops changing, or for 100 generations at most, then prints the generation at
which it stopped and the final board. (\textit{Hint: \token{==} and \token{!=}
  compare two slices of slices element by element.})

\subsection*{Exercise 8}

The \textit{median} of a list of numbers is the number in the middle once the
list is sorted: half of the others are below it, half above. Write a function
\token{median(values: [i32])-> i32} that computes the median of a slice of an
odd length, using the function \token{sort} of Listing~\ref{lst:mutability:sort}.
The slice of the caller must not change. Where does the copy go, and how does
that fit the advice of Section~\ref{sec:mutability:copies}?

\subsection*{Exercise 9}

An interval is a pair of integers, its lowest and its highest value. Write a
function \token{widen(ref dmut bounds: (i32, i32), value: i32)} that widens the
interval of its caller just enough to hold \token{value}, with a \token{match}
whose patterns take references to the fields. Use it to find the lowest and the
highest of a slice of readings, starting from an interval that holds only the
first reading.

\vfill%
\pagebreak

\forceevenpage{}

\section*{Solutions}
\markright{\MakeUppercase{Solutions}}%
\subsection*{Exercise 1}

\begin{lstlisting}[style=bashVerb]
[7, 8, 30] [10, 2, 3] [10, 2, 3] [1, 20, 3]
\end{lstlisting}

\token{b} and \token{c} share the elements of \token{a}, and \token{d} has
elements of its own. \token{c[0] = 10;} writes into the shared elements, which
\token{b} sees too, and \token{d[1] = 20;} changes only \token{d}. Then
\token{a} is given a new slice: from there, \token{a[2] = 30;} changes only the
new elements, and \token{b} and \token{c} keep the first ones.

\subsection*{Exercise 2}

\token{reset} takes an array, which borrows nothing, so a \token{dmut}
parameter is refused: it would change a copy. The function takes a reference
to the array of its caller, \token{ref dmut}, and the call passes
\token{ref counts}. The loop of \token{scale} declares a mutable copy of each
element, which is refused, and a reference iterator cannot go through a
parameter (Section~\ref{sec:mutability:loops}): the loop goes through the
indices. Last, \token{scale} changes the elements of \token{prices}, so the call
passes \token{alias prices}.

\begin{lstlisting}[style=coloredverbatim, escapechar=@, caption=Solution for exercise 2]
use std::io;

fn reset(@\hcb{ref dmut}@ values: [i32 ; 3]) {
  for i in 0 .. 3 {
    values[i] = 0;
  }
}

fn scale(dmut values: [i32], factor: i32) {
  @\hcb{for i in 0 .. values.len \{}@
    @\hcb{values[i] = values[i] * factor;}@
  }
}

fn main() {
  let dmut counts = [4, 5, 6];
  reset(@\hcb{ref}@ counts);
  let dmut prices = copy [10, 20];
  scale(@\hcb{alias}@ prices, 3);
  println(counts, " ", prices);
}
\end{lstlisting}
\vspace{-10pt}%
\begin{lstlisting}[style=bashVerb]
[0, 0, 0] [30, 60]
\end{lstlisting}

\subsection*{Exercise 3}

The loop exchanges the element at index \token{i} with the one at the same
distance from the end, and goes only halfway, since the second half would
exchange the same pairs back. With five elements, the middle one stays in place.

\lstinputlisting[style=coloredverbatim, caption={Solution for exercise 3, \textit{examples/mutability/solutions/reverse.yr}}]{examples/mutability/solutions/reverse.yr}
\vspace{-10pt}%
\begin{lstlisting}[style=bashVerb]
[5, 4, 3, 2, 1]
[4, 3, 2, 1]
[]
\end{lstlisting}

\subsection*{Exercise 4}

The function builds a new slice of the elements it keeps, and gives it to the
variable of its caller through the reference.

\lstinputlisting[style=coloredverbatim, caption={Solution for exercise 4, \textit{examples/mutability/solutions/remove.yr}}]{examples/mutability/solutions/remove.yr}
\vspace{-10pt}%
\begin{lstlisting}[style=bashVerb]
[3, 4, 5]
[3, 0, 4, 0, 0, 5]
\end{lstlisting}

\token{before} does not change. It shares the elements that \token{readings}
held before the call, and \token{removeAll} does not change them: it gives
\token{readings} a new slice, as Section~\ref{sec:mutability:sharing} gave
\token{primes} one.

\subsection*{Exercise 5}

The stock changes in place, so \token{deliver} takes it as a \token{dmut}
parameter, and the call passes it with \token{alias}. The delivery is only read.

\lstinputlisting[style=coloredverbatim, caption={Solution for exercise 5, \textit{examples/mutability/solutions/stock.yr}}]{examples/mutability/solutions/stock.yr}
\vspace{-10pt}%
\begin{lstlisting}[style=bashVerb]
apples: 12
pears: 10
plums: 20
\end{lstlisting}

\subsection*{Exercise 6}

Both programs print the same wrong generations. Generation 1 is:

\begin{lstlisting}[style=bashVerb]
......
..#...
.##...
......
......
\end{lstlisting}

With \token{copy}, \token{before} has a block of rows of its own, but each of
its rows designates the cells of a row of \token{board}
(Figure~\ref{fig:mutability:copies}). With nothing, \token{before} shares
everything with \token{board}. In both cases, when the loop writes the new state
of a cell, \token{before} sees it, and the cells that follow count their
neighbours on a mix of the two generations. Only \token{dcopy} gives
\token{before} cells of its own.

\subsection*{Exercise 7}

\token{step} compares the new board with \token{before} once the loop is done.
The loop of \token{main} counts the generations that changed the board: the
row of four becomes a \textit{beehive}, a shape that no longer changes, at
generation 2.

\lstinputlisting[style=coloredverbatim, caption={Solution for exercise 7, \textit{examples/mutability/solutions/life\_stable.yr}}]{examples/mutability/solutions/life_stable.yr}
\vspace{-10pt}%
\begin{lstlisting}[style=bashVerb]
Stable from generation 2:
.......
.......
..##...
.#..#..
..##...
.......
.......
\end{lstlisting}

\subsection*{Exercise 8}

\token{median} only reads its argument, so its parameter is a plain slice, and
\token{sort} cannot sort it in place: the function copies it, and sorts the
copy. The advice of Section~\ref{sec:mutability:copies} still holds for
\token{sort}, which changes elements and leaves the copy to its caller.
\token{median} is such a caller: it needs the original to stay as it is, so it
makes the copy.

\lstinputlisting[style=coloredverbatim, caption={Solution for exercise 8, \textit{examples/mutability/solutions/median.yr}}]{examples/mutability/solutions/median.yr}
\vspace{-10pt}%
\begin{lstlisting}[style=bashVerb]
Median: 19
Times:  [42, 7, 19, 3, 25]
\end{lstlisting}

\subsection*{Exercise 9}

Each arm takes a reference to the field it may change, and its guard reads the
same field. The interval is a variable of \token{main}, passed as \token{ref
  bounds}; inside \token{widen}, \token{bounds} is a reference, and a variable
that \token{match ref} accepts.

\lstinputlisting[style=coloredverbatim, caption={Solution for exercise 9, \textit{examples/mutability/solutions/bounds.yr}}]{examples/mutability/solutions/bounds.yr}
\vspace{-10pt}%
\begin{lstlisting}[style=bashVerb]
From 6 to 21
\end{lstlisting}

